\begin{abstract}
  Sistemas reator-separador com corrente de reciclo apresentam elevada relevância industrial, pois a recirculação de reagentes pode favorecer o aumento da conversão global e da seletividade do processo.
  Entretanto, o forte acoplamento entre reação, separação e reciclo torna sua dinâmica mais sensível às condições operacionais, exigindo uma análise cuidadosa de seu comportamento.
  Este trabalho caracteriza a dinâmica de um sistema composto por dois reatores contínuos de tanque agitado e um separador, integrados por uma corrente de reciclo, empregando ferramentas clássicas de análise de sistemas lineares nos domínios do tempo contínuo, da frequência e do tempo discreto.
  O modelo fenomenológico não linear foi simulado em malha aberta para identificar qualitativamente a influência de cada entrada sobre os estados do processo.
  Em seguida, a partir do modelo linearizado, obteve-se a matriz de funções de transferência, que permitiu verificar a estabilidade do sistema, determinar a solução analítica da resposta temporal e identificar suas principais características dinâmicas.
  No domínio da frequência, o diagrama de Bode permitiu analisar a resposta do sistema a diferentes frequências, enquanto a Transformada de Fourier fundamentou o desenvolvimento de filtros passa-baixa para atenuação de ruído.
  Já a discretização, com período de amostragem definido com base no teorema de Shannon, preservou as principais características transitórias da resposta contínua, enquanto o mapeamento dos polos para o plano z confirmou a estabilidade da representação discreta.
  Os resultados demonstram que as ferramentas de análise de sistemas lineares permitiram caracterizar a dinâmica do sistema reator-separador com reciclo nos domínios do tempo, da frequência e do tempo discreto.
  Conclui-se que essa caracterização amplia a compreensão do comportamento dinâmico do processo.

  \noindent
  \textbf{Palavras-chave}: reator-separador; reciclo; dinâmica de processos; estabilidade.
\end{abstract}
